\section{The bounded-order characters at exponent one}
\label{sv-section}

The part of case $V$ with small $2$-power conductor admits a
direct use of the coarse Selmer bound. The relevant distinction
is that the degree of the field killing its character is bounded
in terms of the fixed data. The growing conductor therefore
occurs only in the discriminant, and the coarse bound becomes a
fixed power of that conductor. We prove this with the exact
threshold used for case $V_s$ in the submitted source.
\cite{alpo-uniform-source}

Retain $\chi,M,o,E_h,N_h$ from (IV1), and fix an odd prime $p$.
Let $\theta$ be a character of $\operatorname{Gal}(E_h/\mathbb Q)$.
Write $2^c$ for the primitive conductor of its local component
$\theta_2$. Let $\vartheta$ be its part of order prime to $p$,
and put
\[
 \xi=\vartheta/\vartheta_2,\quad
 u=\operatorname{ord}(\xi),\quad
 f=\operatorname{ord}_{\operatorname{lcm}(4,u)}(p),\quad
 \nu=v_2(p^f-1),\quad
 R=16o(pM)^2\,4^\nu.
                                                               \tag{SV1}
\]
The multiplicative order defining $f$ is well defined: $u$ is
prime to $p$ and $p$ is odd. All character orders are finite.

\subsection{The field degree is bounded before the conductor grows}

\begin{lemma}\label{sv-order}
The order of the product of all odd-prime components of
$\theta$ divides $2o$. Consequently $u\mid2o$ and
\[
 \begin{split}
 0<f&\le\varphi(4o),\qquad 2^\nu<p^{\varphi(4o)},\\
 R&\le H_p:=16o(pM)^2p^{2\varphi(4o)}.
 \end{split}                                                   \tag{SV2}
\]
If $2^c<R$, then
\[
                       \operatorname{ord}(\theta)\le2oH_p.
                                                               \tag{SV3}
\]
\end{lemma}
\begin{proof}
At an odd prime the inertia group of $E_h/\mathbb Q$ injects
into the product of the inertia groups of
$K_\chi/\mathbb Q$, $\mathbb Q(\sqrt{-q})/\mathbb Q$, and
$\mathbb Q(\sqrt{-d})/\mathbb Q$. The field $R_h$ is
unramified there. The first group has exponent dividing $o$;
the other two have exponent dividing two. Thus each odd local
character has order dividing $2o$. The order of their product
is the least common multiple of their orders, because their
arguments lie in different Chinese remainder factors. It too
divides $2o$. Taking the part of order prime to $p$ gives
$u\mid2o$.

Set $m=\operatorname{lcm}(4,u)$. Then $m\mid4o$ and
$(m,p)=1$. The order $f$ divides $\varphi(m)$ by Lagrange's
theorem. For positive integers $m\mid n$, one has
$\varphi(m)\mid\varphi(n)$: this follows at once from
$\varphi(r^a)=r^{a-1}(r-1)$ by increasing each prime exponent
and then adding the missing primes. Hence
$f\le\varphi(4o)$. The positive integer $p^f-1$ is divisible
by $2^\nu$, so
\[
                 2^\nu\le p^f-1<p^f\le p^{\varphi(4o)}.
\]
Squaring proves the bound for $R$.

The order of $\theta_2$ divides $\varphi(2^c)$ and hence
divides $2^c$; this also holds for $c=0$. The order of
$\theta$ consequently divides $2o\,2^c$. The hypothesis
$2^c<R\le H_p$ implies (SV3).
\end{proof}

\begin{lemma}\label{sv-field}
Let $F_\theta$ be the fixed field of $\ker\theta$ in
$\mathbb Q(\mu_{N_h})$, and put $F=F_\theta(\mu_p)$.
Under $2^c<R$ this is an abelian, totally imaginary field with
\[
 \begin{split}
 \mu_p&\subset F,\qquad \theta|_{G_F}=1,\qquad
 F\subset\mathbb Q(\mu_{pN_h}),\\
 n=[F:\mathbb Q]&\le n_0:=2opH_p,\qquad
                   |\operatorname{disc}F|\le(pN_h)^n.
 \end{split}                                                   \tag{SV4}
\]
Here $n_0$ depends only on $p,M,o$.
\end{lemma}
\begin{proof}
The first statements follow from the definitions. The image of
$\theta$ is cyclic of order $\operatorname{ord}(\theta)$, so
$[F_\theta:\mathbb Q]=\operatorname{ord}(\theta)$.
Adjoining $\mu_p$ increases this degree by at most $p-1$.
Now use (SV3). The presence of $\mu_p$, with $p$ odd, makes
$F$ totally imaginary.

For completeness, if $F\subset L=\mathbb Q(\mu_m)$ and
$n=[F:\mathbb Q]$, the discriminant tower formula implies
\[
             |\operatorname{disc}F|^{[L:F]}
                                      \le|\operatorname{disc}L|.
\]
Moreover $|\operatorname{disc}L|\le m^{[L:\mathbb Q]}$.
One way to obtain the latter inequality without the explicit
cyclotomic discriminant formula is to write
$X^m-1=\Phi_m(X)Q(X)$. At a primitive root $\zeta$,
\[
                 m\zeta^{m-1}=\Phi'_m(\zeta)Q(\zeta).
\]
Taking absolute norms shows
$|N_{L/\mathbb Q}\Phi'_m(\zeta)|\le m^{[L:\mathbb Q]}$,
because $Q(\zeta)$ is a nonzero algebraic integer. The
field discriminant divides the discriminant of the order
$\mathbb Z[\zeta]$, whose absolute value is this derivative
norm. Combining the two inequalities gives
$|\operatorname{disc}F|\le m^n$. Apply this with $m=pN_h$.
\end{proof}

\subsection{Restriction with a bounded invariant module}

Use the normalized length $s_p(1,\theta)$ of (IV6), with an
arbitrary finite coefficient extension containing the values of
$\theta$. Write $b$ for its ramification index.

\begin{lemma}\label{sv-restriction}
For the field $F$ in (SV4),
\[
                     s_p(1,\theta)\le\lambda_p(F,1)+n^2.
                                                               \tag{SV5}
\]
In particular, finiteness of the group on the right implies
finiteness on the left.
\end{lemma}
\begin{proof}
Since $F$ is finite and contains $\mu_p$, its intersection
with $\mathbb Q(\mu_{p^\infty})$ is
$\mathbb Q(\mu_{p^a})$ for an integer $a\ge1$.
Indeed the Galois group above $\mathbb Q(\mu_p)$ is the
procyclic group $1+p\mathbb Z_p$, whose open subgroups are
$1+p^a\mathbb Z_p$. The degree of this intersection gives
\[
                            (p-1)p^{a-1}\le n,
\]
and in particular $a\le n$. On $G_F$, the coefficient
module is $(K/\mathcal O)(-1)$. The cyclotomic image is
$1+p^a\mathbb Z_p$, so its invariant subgroup is precisely
$(K/\mathcal O)[p^a]$, of length $ba$.

Apply the inflation--restriction sequence (IV9), with this
$F$ and coefficient module. The kernel has length at most
$n\,ba$, since its cocycles form a submodule of the functions
from the group of order $n$ to that invariant subgroup.
The image belongs to the strict Selmer group over $F$,
whose coefficient-extended length is $b\lambda_p(F,1)$
by (IV10). Division by $b$ proves (SV5), since $a\le n$.
The argument again requires no coprimality between $p$ and $n$.
\end{proof}

\begin{theorem}\label{sv-exact-case}
For fixed $p,M,o$, there is an integer $C$ such that, for every
$h,q,d,\theta$ satisfying the preceding conditions and $2^c<R$,
the strict Selmer group at $e=1$ is finite and
\[
                           s_p(1,\theta)\log p\le(pN_h)^C.
                                                               \tag{SV6}
\]
The same conclusion holds with the source's extra zero local
conditions at the primes dividing $N_h$.
The exponent $C$ is independent of $h,q,d$ and of the coefficient
extension. Thus this proves the source's case $V_s$ for every
$p$ in its fixed finite set of odd primes.
\end{theorem}
\begin{proof}
The field in (SV4) is totally imaginary and contains $\mu_p$,
so the strict Selmer bound of Theorem~\ref{kt-fourteen}, applied
with $E=1$, proves finiteness and
\[
 \lambda_p(F,1)\log p
 \le\exp\bigl((2n)^{14}\log(2|\operatorname{disc}F|)\bigr).
                                                               \tag{SV7}
\]
Put $x=pN_h\ge3$. By (SV4),
$2|\operatorname{disc}F|\le x^{n+1}$. Hence the right side
of (SV7) is at most $x^A$, where
\[
                            A=(2n_0)^{14}(n_0+1)
                                                               \tag{SV8}
\]
is fixed. By (SV5),
\[
                       s_p(1,\theta)\log p
                                      \le x^A+n_0^2\log p.
\]
Choose an integer $k\ge0$ with $3^k\ge n_0^2\log p$.
Then the latter expression is at most
$2x^{\max(A,k)}\le x^{\max(A,k)+1}$, proving (SV6).
Additional local conditions decrease the group.

This proof uses the algebraic $K$-theory bound directly;
the finite-field homology and arithmetic-cycle comparison
is not needed. The numerical size of the exponent is immaterial here:
$n_0$ is fixed before $h,q,d$ vary. Applying the same argument to a field
whose degree grew with $h$ would not give a fixed exponent,
and no such step is used.
\end{proof}

\begin{corollary}\label{sv-summed}
For a fixed finite set $P$ of odd primes, there is an integer
$C_P$, depending only on $P,M,o$, such that for every $t\ge2$
\[
 \sum_{p\in P}\ \sum_{\substack{\theta\in
          \widehat{\operatorname{Gal}(E_h/\mathbb Q)}\\2^{c(\theta)}<R_p(\theta)}}
                   s_p(1,\theta)\log p\le(N_ht)^{C_P}.
                                                               \tag{SV9}
\]
Thus the entire case $V_s$ contribution has a fixed polynomial
bound before selecting $h,q,d$.
\end{corollary}
\begin{proof}
For fixed $p$, (SV2) bounds the $2$-part of the primitive
conductor by $H_p$. If $T$ is the largest power of two at most
$H_p$, all these local components factor modulo $T$, so there
are at most $\varphi(T)\le H_p$ of them. At the odd primes
of $M$ there are at most $\varphi(M/2^{v_2(M)})\le M$
choices; at each of $q,d$ there are at most two. Thus there
are at most $4MH_p$ characters in this sum, independently of
$h,q,d$.

Choose $C_p$ from Theorem~\ref{sv-exact-case}. The sum is at
most $A N_h^{C_0}$, where
\[
 C_0=\max_{p\in P}C_p,\qquad
 A=\sum_{p\in P}4MH_p p^{C_p}.
\]
For nonempty $P$, choose $r\ge0$ with $2^r\ge A$.
Then $A N_h^{C_0}\le(N_ht)^{C_0+r}$ for $t\ge2$.
For empty $P$ the sum is zero. All choices involve only the
fixed data.
\end{proof}


